\begin{table}[H]
\centering
\caption[Citywide Simulated Resistance Penalties]{\textbf{Citywide Simulated Resistance Penalties by Planning Scenario.} Expected processing delay penalty and hurdle risk (probability of petition-induced withdrawal) if subjected to a 20\% statutory neighborhood petition, simulated across all $N=271,567$ addressable Austin parcels. The scenarios represent common entitlement requests: a strict use-change with no physical expansion (+0 ft), a "missing middle" medium-density upzone (+25 ft), and a mid-rise transit corridor rezoning (+55 ft).}
\label{tab:citywide_simulation}
\begin{tabular}{l cc cc}
\toprule
& \multicolumn{2}{c}{\textbf{Delay Penalty (Days)}} & \multicolumn{2}{c}{\textbf{Hurdle Risk (\%)}} \\
\cmidrule(lr){2-3} \cmidrule(lr){4-5}
\textbf{Scenario} & \textbf{Mean} & \textbf{Median} & \textbf{Mean} & \textbf{Median} \\
\midrule
Use-Change (+0 ft) & 568 & 99 & 22.7\% & 0.0\% \\
Missing Middle (+25 ft) & 735 & 119 & 25.3\% & 0.0\% \\
Transit Corridor (+55 ft) & 740 & 121 & 24.9\% & 0.0\% \\
\bottomrule
\multicolumn{5}{l}{\footnotesize Generated via continuous spatial causal inference across all municipal parcels.}\\
\end{tabular}
\end{table}
